% res_join.tex -- Results 4.3: joining the team (RQ3). 17 Oct (v3_full): sim_attach.tex and skel_join.tex merged.
% The join tables fly on a second code tree after the 21-22 Oct rig day (Wesley 16 Oct, inventory decisions (1), (2),
% (9), (11)): T11 and F14 re-flown on the co-motion confirmation; T12 two hover runs under a card, the circle
% exploratory; the push-abort clearance bar is the newcomer's own clearance. The 12 Oct contact-and-lift T11 and F14 are
% STALE-MECH (one line of history in App D). Tonight's exploratory join re-fly (17-18 Oct) is not thesis evidence.
% Rig record: docs/reports/supervisor_brief_2026-10-16.md and morning_report_2026-10-16.md (rig bags re-read), registry
% RIG-1014-t6h_4, RIG-1015-a1, a2_2, a3_2; the offline replay: card docs/experiments/2026-10-15_attach_comotion_rig.md
% (baseline; tools/proof_replay.py). Fault coverage: drop_takes_back card (16 Oct, twin, SUPPORTED WITH CORRECTIONS),
% quad4_tilt_cut rev 4 (week5 R1702-R1705, 4 of 4, exploratory), land_every_phase (9 Oct), newcomer_land_abort.

\subsection{Joining the team}
\label{sec:res_join}

This section answers RQ3: what must be observed before a UAV flown by another controller is admitted to a carried
ring, and what does admission cost the ring? It is the least complete of the three. The confirmation was rewritten on
14 October, after the rig's first catch was rejected by the proof of the day, and the twin claim for the new
confirmation flies on the join code tree after the rig day of 21--22 October; its tables and figures below are
placeholders sized for that batch. What is in hand is the contract itself (Section~\ref{sec:admission}), an offline
replay of the new confirmation on the rig's one catch, the rig's attach record so far, and the fall-backs.

\subsubsection{The admission contract in the twin}
\label{sec:res_contract}

Figure~\ref{fig:F14} and Table~\ref{tab:T11} will give the contract on the join code tree: the collaborator's
controller flies the newcomer down onto plate 9 of a carried ring, and this work confirms the catch from motion
capture, switches the radio and folds the newcomer in, or takes it back. Smoke runs of the co-motion confirmation in
the twin admitted the newcomer, folded it in and landed (exploratory, one run per code tree); the twin's rigid tether
cannot go slack, so the rig's thresholds stay on the rig side.

\pendingfig{fig:F14}{The co-motion confirmation in the twin: a newcomer flown by the collaborator's controller comes
down onto plate 9 of a hovering ring and the catch is confirmed from motion capture alone. (a) Heights of the
newcomer's body, its magnet tip and the plate point; (b) the tip's drift from its contact point in the plate frame,
against the 6\,mm confirmation and the 45\,mm loss; (c) the newcomer's throttle across the switch of its radio; (d) the
ring's tilt departure against the 5$^\circ$ bar.}{the hover attach runs of the join claim, on the join code tree}

\begin{table}[htbp]
  \centering
  \footnotesize
  \caption[The twin claim for the join]{The twin claim for the join (T11): the collaborator's controller flies the
  newcomer to contact and this work's co-motion confirmation admits it. Ranges over $n$ runs on one code tree; the
  ring's tilt rise on the release metric with the contact as the event; clearance: the newcomer's body and rod to every
  carrier, and in the push abort the newcomer's own. [Every cell: the join claim on the join code tree, after the
  21--22 October rig day, then the reviewer.]}
  \label{tab:T11}
  \begin{tabularx}{\textwidth}{XcYYYYl}
    \toprule
    Claim line & $n$ & Contact to confirmation (s) & Throttle step (\%) & Ring tilt rise ($^\circ$) & Clearance (m)
      & Reviewer \\
    \midrule
    Attach, hover & [\dots] & [\dots] & [\dots] & [\dots] & [\dots] & [\dots] \\
    Take-back, hover & [\dots] & -- & -- & [\dots] & [\dots] & [\dots] \\
    No catch, hover & [\dots] & -- & -- & [\dots] & [\dots] & [\dots] \\
    LAND during the confirmation & [\dots] & -- & -- & [\dots] & [\dots] & [\dots] \\
    No catch, circle & [\dots] & -- & -- & [\dots] & [\dots] & [\dots] \\
    LAND take-back, circle & [\dots] & -- & -- & [\dots] & [\dots] & [\dots] \\
    Push abort & [\dots] & -- & -- & [\dots] & [\dots] & [\dots] \\
    Re-attach & [\dots] & [\dots] & [\dots] & [\dots] & [\dots] & [\dots] \\
    \bottomrule
  \end{tabularx}
\end{table}

\paragraph{Counterfactuals.} [On the join code tree: the newcomer seeded at its full share against 0.1\,N; the
confirmation without its plate-motion condition; the hand-over step against a thrust blend of the kind TRADY uses
\cite{sugihara2023trady}.]

\subsubsection{The confirmation on rig noise}

The twin's magnet catches at the instant of contact and its ring hardly moves, so the confirmation's plate-motion and
taut-tether conditions can only be judged on rig data. Replayed offline over the flights of 14 October with the rig's
thresholds (a 24\,mm catch radius and each flight's own free hang), the co-motion rules confirm the one catch of that
day 2.81\,s after its contact, the tip drifting 1.1\,mm in the plate frame while the plate moved 16\,mm, and find no
contact and no false confirmation in the other descents. [False accepts and rejects over every rig recording with a
descent, and the rig thresholds set from them.]

\subsubsection{On the rig}
\label{sec:res_join_rig}

Eight attach flights on 14 and 15 October, the collaborator's controller flying the newcomer from its mat to plate 9
of a ring held at 1.0\,m by three carriers, gave one contact and no fold-in (exploratory). In that contact the magnet
held plate 9 through three lifts, and the proof of the day rejected it on its taut check; the co-motion confirmation
that replaced it confirms it offline (above). In the other descents the magnet was inside the catch radius at the
bottom only 22\,\% of the time; three bottom holds ended with the tip at plate height 3.3--4.3\,cm along the ring from
the plate point, past the edge of the 60\,mm plate; and the ring bounced about 0.8\,cm at 1.8\,Hz under the newcomer,
which followed the bounce 0.25\,s late. In one flight the collaborator's match check never held long enough for a
descent (42\,s), and in another the newcomer's tether passed 0.10\,m from a carrier on the way in. One of these flights
also lost the ring when a carrier's magnet let go (L11, Section~\ref{sec:res_leave}). The attach on the rig is
therefore limited by getting the magnet onto the plate, which is the collaborator's controller's task and the plate's
size, before any admission is tested. The response is hardware first: 90\,mm plates, inside whose catch radius the
magnet would have been 49\,\% of the time at the bottom on these flights instead of 22\,\%, and 1\,m tethers, which give
the newcomer 0.70\,m of room from the nearest carrier instead of 0.45\,m; both are planned for the next rig day.

\pendingfig{fig:F15}{The first attaches on the rig: the newcomer's descent onto plate 9, the contact, the
confirmation, the switch of its radio and the fold-in, with the ring's tilt against the 5$^\circ$ bar for the thesis
and the 10$^\circ$ safety limit.}{the attach flights of the 21--22 October rig day}

\subsubsection{Fault coverage and authority}
\label{sec:res_join_faults}

Every fall-back of Section~\ref{sec:safety} has been flown in the twin from the phase it guards, the attach
phases on the earlier attach sequence (Appendix~\ref{sec:app_faults}). LAND from every phase of an attach landed and
disarmed every UAV except in one run, LAND 0.5\,s after a weld, where the twin's rigid rod tipped the newcomer (on the code tree of 9 October, before the
collaborator's controller flew the attach). A fault on a newcomer in transit disarms it alone and lands the carriers.
When the ring is dropped while the collaborator's controller flies the newcomer, the newcomer is taken back on the
drop tick, kept at least 0.5\,m from every carrier and set down at its rest 14--17\,s later (twin only); without that
switch the twin reproduced the rig's problem, its newcomer flying on towards the fallen ring. The tilt cut caught an
injected upset on a newcomer under the collaborator's controller and disarmed it in four of four runs (exploratory);
it stays off on the rig. On the rig, a LAND in the collaborator's match phase ended in a take-back with every UAV
landed (15 October).

\subsubsection{The full mission}

\begin{table}[htbp]
  \centering
  \footnotesize
  \caption[The full mission]{The full mission (T12): the UAV on plate 9 is released, picks up and drops the object
  with the collaborator's controller, and rejoins. [Two hover runs under the full mission's card on the join code
  tree; the circle reported as exploratory.]}
  \label{tab:T12}
  \begin{tabularx}{\textwidth}{lcYYYY}
    \toprule
    Path & $n$ & Release tilt rise ($^\circ$) & Reattached & Contact to confirmation (s) & Ring tilt rise at the attach ($^\circ$) \\
    \midrule
    Hover & [\dots] & [\dots] & [\dots] & [\dots] & [\dots] \\
    Circle (exploratory) & [\dots] & [\dots] & [\dots] & [\dots] & [\dots] \\
    \bottomrule
  \end{tabularx}
\end{table}

\paragraph{Answer to RQ3.} [After the join code tree and the rig day: the twin claim, the counterfactuals and the rig
attach.] What can be said now is that what must be observed is defined and observable from motion capture alone: a
contact in the plate's frame and, after the other controller's lift, the tip moving with the plate on a taut tether,
which confirms the rig's one catch offline where the earlier proof rejected it; and that on the rig the attach is
limited before admission, by the magnet reaching the plate. What admission costs the ring is not yet measured on the
current confirmation.
